% Einheit 3 von 3 – Symmetrieebenen von Körpern (bank/spiegelung/e3.jsonl, zusammenbau v0.1)
%% TODO Zeitmarke Einheit 3: Marken-Zeile nicht lesbar
%% TODO Zweigzeile Teil 1: was hier gelernt wird (Satz in Schülersprache aus dem Einheitstitel) – Einheit 3
\einheitenkopf[e3]{Einheit 3 von 3 · Symmetrieebenen von Körpern}
\zweigzeile{Abitur GK}
\setcounter{aufgabe}{14}

%% TODO Titel: Ich-kann-Satz fehlt in der Bank (Platzhalter: Kettenname) – Nr. 15
%% TODO Anweisung: ein Satz über der ersten Teilaufgabe fehlt in der Bank – Nr. 15
\begin{aufgabe}{Symmetrieebenen von Körpern}
%% TODO \rechenplatz{n}: Schrittzahl des Grundfalls fehlt in der Bank (nur bei fester Schreibform, 2.3 b)
\begin{teile}
\teil Welche Koordinate wechselt? $A(3 | -2 | 1)$ und $B(3 | 2 | 1)$ \\ \kreuz{$x_1$ – symmetrisch zur $x_2x_3$-Ebene} \\ \kreuz{$x_2$ – symmetrisch zur $x_1x_3$-Ebene} \\ \kreuz{$x_3$ – symmetrisch zur $x_1x_2$-Ebene}
\teil Begründe, dass die $x_1x_3$-Ebene Symmetrieebene des Körpers mit den Ecken $A(1 | -3 | 0)$, $B(1 | 3 | 0)$, $C(4 | -3 | 0)$, $D(4 | 3 | 0)$, $E(2 | -1 | 4)$, $F(2 | 1 | 4)$ ist.
\teil Ein Turmdach hat die Ecken $A(0 | 0 | 0)$, $B(10 | 0 | 0)$, $C(10 | 10 | 0)$, $D(0 | 10 | 0)$, die Giebelspitzen $E(5 | 0 | 4)$, $F(10 | 5 | 4)$, $G(5 | 10 | 4)$, $H(0 | 5 | 4)$ und die Spitze $S(5 | 5 | 9)$ (1 LE = 1 m). Genau eine der Ebenen $M_1: x_1 = 10$, $M_2: x_1 + x_2 = 10$, $M_3: x_3 = 4$ ist Symmetrieebene des Dachs. Gib sie an und beschreibe ihre Lage. \hfill (Abitur 2022 GK)
\teil Die Abbildung zeigt den Quader mit $0 \le x_1 \le 4$, $0 \le x_2 \le 2$, $0 \le x_3 \le 3$. Gib die Symmetrieebene an, die senkrecht auf der $x_1$-Achse steht, und zeichne ihre Schnittfigur mit dem Quader ein.

\begin{ksys3}\rquader{0,0,0}{4}{2}{3}\end{ksys3}
\teil Ein Quader hat die Grundfläche $A(1 | 3 | 0)$, $B(7 | 3 | 0)$, $C(7 | 7 | 0)$, $D(1 | 7 | 0)$ und die Höhe $4$. Gib die beiden senkrechten Symmetrieebenen an.
\teil Ein gerades Prisma hat die Grundfläche $A(1 | 0 | 0)$, $B(5 | 0 | 0)$, $C(3 | 4 | 0)$ und die Höhe $4$. Begründe, dass die Ebene $x_1 = 3$ Symmetrieebene ist.
\teil Für $t > 0$ hat die Pyramide $A_tB_tCDS_t$ die Ecken $A_t(3 | 5 + 2t | 0)$, $B_t(-1 | 5 + 2t | 0)$, $C(-1 | -1 | 0)$, $D(3 | -1 | 0)$ und $S_t(1 | t + 2 | 5)$. Gib für jeden Wert von $t$ die beiden Symmetrieebenen an und begründe. \hfill (Abitur 2026 LK)
\end{teile}
\end{aufgabe}

%% TODO Titel: Ich-kann-Satz fehlt in der Bank (Platzhalter: Kettenname) – Nr. 16
%% TODO Anweisung: ein Satz über der ersten Teilaufgabe fehlt in der Bank – Nr. 16
\begin{aufgabe}{Symmetrieebenen von Körpern – Fehler finden, Begründen}
\begin{teile}
\teil Finde den Fehler. Pyramide $ABCDS$ mit $A(4 | 4 | 0)$, $B(-4 | 4 | 0)$, $C(-4 | -4 | 0)$, $D(4 | -4 | 0)$, $S(0 | 0 | 6)$. Behauptung: „Die Ebene $x_3 = 3$ halbiert die Höhe, also ist sie Symmetrieebene.“
\teil Warum genügt eine einzige Punktprobe, um eine Ebene als Symmetrieebene eines Körpers auszuschließen?
\end{teile}
\end{aufgabe}
